
%(BEGIN_QUESTION)
% Copyright 2012, Tony R. Kuphaldt, released under the Creative Commons Attribution License (v 1.0)
% This means you may do almost anything with this work of mine, so long as you give me proper credit

Anta at vi begynner med dette matematiske utsagnet:

$$3 = 3$$

Ikke spesielt oppsiktsvekkende, men likevel matematisk helt sant. Hvis jeg legger til tallet ``5'' på venstre side av ligningen, vil uttrykkene på hver side av likhetstegnet ikke lenger være like. Da må likhetssymbolet erstattes med et ``ikke lik''-symbol:

$$3 + 5 \not = 3$$ 

Hvis vi beholder ``$3 + 5$'' på venstre side, hva må da gjøres på høyre side for å få en likhet igjen?

\underbar{file i01303}
%(END_QUESTION)





%(BEGIN_ANSWER)

Uten å endre venstre side av uttrykket, er den eneste måten å få likhet igjen å legge til tallet ``5'' på høyre side også:

$$3 + 5 = 3 + 5$$ 

%(END_ANSWER)





%(BEGIN_NOTES)

Et grunnleggende prinsipp i algebra er at enhver operasjon er tillatt i en ligning, så lenge den samme operasjonen utføres på begge sider. Denne oppgaven introduserer prinsippet på en skånsom måte, med tallverdier i stedet for variabler.

%INDEX% Mathematics review: basic principles of algebra

%(END_NOTES)
